\section{How the measurements were made}
\label{sec:method}

\paragraph*{Machine.} % Generated by tools/gen_tables.py; do not edit by hand.
An 12th Gen Intel Core i5-12450H laptop processor; Rust 1.92.0, Node.js 24.10.0, gcc 13.3.0 for NBIS.
 The processor has four performance cores
(two threads each) and four efficiency cores; it has AVX2 and FMA but not AVX-512;
Section~\ref{sec:server} repeats the single-core measurements on a server processor that has it.
Single-thread measurements run pinned to one performance core (\code{taskset -c 2}); the
multi-threaded ones use the performance cores, and the batch measurements all twelve hardware
threads. The GPU is the laptop's NVIDIA GeForce RTX 3050 (6 GB).

\paragraph*{Data.} The 40 images of the NIST WSQ reference data set version 2.0 (rolled prints,
plain impressions and four-finger slaps, from 375$\times$526 to 1600$\times$1500 pixels). For
decoding, the inputs are the streams NBIS writes for those images at 0.75 and at 2.25 bits per
pixel, 80 files in all; the 86 files NIST distributes for decoder testing were used for the
correctness checks but not for timing, because their Huffman tables make NBIS print a warning
per table, which would add console output to its times.

\paragraph*{Reference.} NBIS 5.0.0, compiled with its own build system (\code{gcc -O2}), timed
in memory through \code{wsq\_decode\_mem} and \code{wsq\_encode\_mem} by a small C harness
(\code{bench/nbis}). Our codec is timed in memory as well (\code{wsq bench}); file reading and
writing are excluded on both sides.

\paragraph*{Statistic.} Every time is the median of 9 runs (5 for encoding) after one warm-up
run; the decoder and encoder reuse their buffers across runs, as a program decoding many images
would. Per-stage times come from timestamps inside the codec.

\paragraph*{What \emph{faster} means here.} The rows of Tables~\ref{tab:decode}
and~\ref{tab:encode} add up the per-image medians over the 40 images of a set; speedups are
ratios of those sums, so large images weigh more than small ones, as they do in practice.
